% Measurement hazards, standalone v2 paper.
% Numbers from the A2 acceptance gate, analysis_v2, the RETRACTED split sidecars,
% and the H6 degeneracy analysis.

\section{Measurement hazards}
\label{sec:hazards}

Four choices that look like implementation details changed conclusions in this study. We
report them because each is invisible in a finished results table, and because at least two
of them are made silently by default in commonly used tooling.

\subsection{Numerical precision moves roughly one label in ten}

We ran the acceptance gate for this study at three precisions on the same model, comparing
episode-level verdicts against \texttt{float32}. Both \texttt{bfloat16} and
\texttt{float16} fell below the pre-registered 95\% agreement bar, at 88.4\% and 89.6\%
respectively. Roughly one episode in ten receives a different verdict from the arithmetic
alone, with no change to the model, the prompt, or the decoding strategy.

This is larger than several of the effects reported in this paper. A flip-rate difference of
five percentage points between two conditions is a real finding under one precision and
potentially an artifact under another. We therefore ran everything in \texttt{float32},
which is expensive and which caps the model sizes reachable on a single accelerator, and we
would encourage anyone reporting small behavioural differences between conditions to state
their precision and check it.

\subsection{Conflating flips with abandonment can reverse a comparison}

When a model stops defending a correct answer, it can either commit to a wrong one or commit
to nothing. Our earlier work argued these are different failure modes and should be counted
separately. The stronger version of that claim is available here: conflating them can flip
the sign of a headline comparison.

For Qwen-1.5B, the emotional versus simple comparison gives an odds ratio of 0.41 in favour
of simple pushback when restricted to flips, at $p = 0.003$. Counting abandonment as
destabilisation gives 1.47 in the opposite direction, at $p = 0.079$. The same episodes,
the same test, opposite conclusions. Qwen-7B is stable under both codings, 0.31 against
0.53, and Llama-8B strengthens, 2.37 against 4.10. So the hazard is not universal, but it is
not detectable from the aggregate either.

\subsection{Paraphrases that omit a reconsideration cue provoke premise refusal}

Our five paraphrases per pushback type were written to vary wording while holding pressure
type fixed. They do not entirely succeed. Paraphrase 0, the sentence used in our earlier
work, produces 0.9\% abandonment; paraphrases 1, 3 and 4 produce 25\% to 30\%, concentrated in the social and emotional types. The localisation is model-dependent: in Llama-1B the deflections fall mainly on authoritative pushback instead.

Inspection shows why. The high-abandonment paraphrases dropped an explicit request to
reconsider or double-check, and models respond to the framing of the pushback instead of the
question. A social paraphrase asserting that the model's friends disagree elicits statements
about not having friends. These are premise refusals, not answer abandonment, and grouping
them together inflates apparent abandonment.

We chose not to amend the frozen templates, since tuning stimuli to observed model behaviour
after seeing data is exactly what pre-registration is meant to prevent. Instead we added a
second grading pass over abandonment verdicts only, separating genuine withdrawal from
premise refusal. It changes no flip counts. The correction is large: between 50\% and 99\%
of raw abandonment verdicts are premise refusals, at 98.6\% for Qwen-0.5B and 50.0\% for
Llama-3B.

Two consequences. First, raw abandonment rates in this literature should be treated as upper
bounds unless the failure mode has been inspected. Second, our earlier cross-family claim
about abandonment survives the correction: genuine abandonment counts remain roughly an
order of magnitude higher in Llama than in Qwen at comparable scale, so that finding was not
an artifact of the templates even though the raw numbers were inflated by them.

\subsection{Hooked generation degenerates on one model family}
\label{sec:degeneracy}
Causal ablation requires generating with hooks installed, which in practice means a
different inference implementation from the batched engine used for behavioural runs. On
Llama-3.1-8B, greedy \texttt{float32} generation through the hooked implementation
degenerates into repetition loops on 13\% of episodes, by a 5-gram repetition test, where
the batched engine on identical weights produces under 1\%. Qwen-7B shows no such effect
under the same code path, at 2.0\% and 1.5\% abandonment in the two arms.

Because the degeneracy is symmetric across the ablated and baseline arms, 94 episodes each,
it does not bias the contrast in Table~\ref{tab:h6}, and the Llama-8B null survives both
codings of the affected episodes. It does limit external validity: that causal test was run
on a model behaving measurably differently from the one in our behavioural results.

The obvious remedy, a repetition penalty, is unavailable in this configuration, because the
implementation applies frequency penalties only inside its sampling branch and greedy
decoding at temperature 0 bypasses it. Fixing this properly means either a custom greedy
loop with $n$-gram blocking or a move to stochastic decoding, and either would require
rerunning both models for comparability. We flag it for anyone building ablation studies on
this stack, since a silent 13\% degeneracy rate in one arm of a paired comparison would be
easy to miss.
 
\subsection{Parametric thresholds fail on heavy-tailed head-level nulls}
\label{sec:pointwise}
 
The head-level gate in Section~\ref{sec:gate} carries two null terms: a pointwise threshold,
the mean plus two standard deviations of that head's own shuffled-label distribution, and a
max-statistic threshold, the 95th percentile of per-permutation maxima across every head.
The first is inherited from the residual-stream gate. It does not survive the move to head
level.
 
The symptom is visible in Table~\ref{tab:h2}. Qwen-7B's best head reaches AUROC 0.840 against
a max-null of 0.821, and twelve heads clear that corrected threshold, yet the model is
recorded as failing because the best head does not clear its own pointwise threshold of
0.924. A pointwise threshold exceeding the corrected family-wise threshold is
arithmetically odd: the maximum over 784 positions should dominate any single position's
mean plus two standard deviations. In Qwen-7B it does not for 150 of the 784 heads.
 
The cause is that per-head null distributions are heavy-tailed, so a mean and a standard
deviation are a poor summary of them. The median pointwise threshold across Qwen-7B's heads
is 0.727, implying a null standard deviation near 0.11, against 0.607 and roughly 0.054 for
Qwen-3B. Where the null is that dispersed, adding two standard deviations produces a
threshold no observed value can reach, and the term rejects real signal rather than
controlling false positives. We checked and rejected the obvious explanation that these are
degenerate heads with near-zero activations: Qwen-7B's best head has an activation norm of
5.98 against a median of 2.75 across heads, so it is among the stronger heads, not a dead
one.
 
We report the pre-registered gate as primary because it was fixed in advance, and we report
the max-statistic result alongside it as a sensitivity analysis. Under the max-statistic
term alone, Qwen-7B and Llama-8B pass at head level and Qwen-3B passes marginally. This does
not change the conclusion of Section~\ref{sec:h2}, since under either reading head-level
decodability tracks residual-stream strength rather than compensating for it, and the models
with clean residual nulls have no surviving heads under either term. It does mean that
anyone extending a residual-stream validation protocol to a scan over hundreds or thousands
of positions should use a non-parametric multiplicity correction and drop the pointwise term,
rather than carrying both across as we did.